\section{Добавление к вопросу о сравнимости относительных топологий}





В связи с использованием в разделе 5 соглашения ${\bold K} \in


{\frak K}$ отметим некоторые свойства, дополняющие (5.3).


Введем в рассмотрение топологию


${{\tau}_\otimes}({\cal L})$ ([5], c.\,80; [6], c.\,44)


поточечной сходимости в


${\Bbb A}({\cal L})$: ${{\tau}_\otimes}({\cal L})$ индуцирована на


${\Bbb A}({\cal L})$ топологией ${{\otimes}^{\cal L}}({\tau}_{\Bbb R})$


множества ${{\Bbb R}^{\cal L}}$, соответствующей тихоновскому


произведению экземпляров ${\Bbb R}$ в обычной $|\cdot|$-топологии


${\tau}_{\Bbb R}$ при использовании ${\cal L}$ в качестве индексного


множества. При $M \in {\cal P}({\Bbb A}({\cal L}))$ семейство


${{\tau}_M^\otimes}({\cal L}) \stac {{\tau}_\otimes}({\cal L})\mid_M$


есть топология, индуцированная ([20], c.\,77) на $M$ из


$({{\Bbb R}^{\cal L}},{{\otimes}^{\cal L}}({\tau}_{\Bbb R}))$. Пусть


([6], c.\,42)


$$


{{\frak C}_*}[{\cal L}] \stac \{S \in {\cal P}({\Bbb A}({\cal L})) \mid


{N_{{{\tau}_S^\otimes}({\cal L})}}(\mu) \cap {{\bold B}_*}({\cal


L}) \ne \emp


\ \fo{\mu} \in S\}.


$$





Используя свойства типа соотношения (3.5.6) из


[6], можно показать, что


([5], c.\,94) $\fo{M} \in {{\frak C}_*}[{\cal L}]$


$$


{{{\tau}_*}({\cal L})\mid_M} = {{\tau}_M^\otimes}({\cal L}).


$$


При $B \in {{\bold B}_*}({\cal L})$ имеем вложение ${(\add)_+}[{\cal L}] \cup


B \in {{\frak C}_*}[{\cal L}]$.





Отметим аналог последнего свойства, связанный с использованием


алгебраических сумм п/м ${\Bbb A}({\cal L})$. Если $A \in


{\cal P}({\Bbb R}^{\cal L})$


и $B \in {\cal P}({\Bbb R}^{\cal L})$, то $A{\oplus}B \stac \{\mu + \nu:


(\mu,\nu) \in A \times B\}$.


В этих обозначениях $\fo{B} \in {{\bold B}_*}({\cal L}):


{(\add)_+}[{\cal L}]{\oplus}B \in {{\frak C}_*}[{\cal L}]$.


При $M \in {{\frak C}_*}[{\cal L}]$ имеем


${\cal P}(M) \subset {{\frak C}_*}[{\cal L}]$; здесь аналогия с


(5.3). Полагая ${{\frak C}'_*}[{\cal L}] \stac {{\frak C}_*}[{\cal


L}] \setminus \{\emp\}$, имеем


${{\frak C}'_*}[{\cal L}] \subset {\frak K}$. С учетом (5.3) получаем, что


(5.2) --- весьма обширное семейство п/м ${\Bbb A}({\cal L})$, и


предположение ${\bold K} \in {\frak K}$ раздела 5 не выглядит


ограничительным.


\medskip





\centerline{\large{\bf Часть II}}


\medskip





Рассматриваются


конструкции расширения абстрактных задач управления,


прототипом которых являются задачи управления линейными системами


с разрывными коэффициентами при управлении. Сам метод построения


расширений является конкретизацией общих процедур первой части


данной работы.





\section{Постановка задачи}





В дальнейшем используются понятия и обозначения первой части.





Рассмотрим одну типичную ситуацию, возникающую при управлении линейными


системами. Пусть дано управляемое дифференциальное уравнение


\begin{equation}


\Dot x (t) = {A(t)}{x(t)} + {f(t)}{b(t)}


\end{equation}


в $n$-мерном фазовом пространстве ${\Bbb R}^n$;


движения (8.1) рассматриваются


на конечном промежутке времени ${I_0} \stac [{t_0},{{\vt}_0}]$,


где ${t_0} \in {\Bbb R}$ и ${{\vt}_0} \in ]{t_0},\infty[$.


В (8.1) $A(\cdot)$ есть $n \times n$-матрицант на $I_0$ с непрерывными


в/з компонентами ${A_{i,j}}(\cdot)$, $i \in \overline{1,n}$, $j \in


\overline{1,n}$.


Через $b = b(\cdot)$ в (8.1) обозначено отображение из


$I \stac [{t_0},{{\vt}_0}[$ в ${{\Bbb R}^n}$. С $I$ естественно


связана полуалгебра (п/а) ${\cal I} \stac \{[u,v[:


(u,v) \in {I_0} \times {I_0}\}$ пространства-стрелки; пространство


${B_0}(I,{\cal I})$ всех ${\cal I}$-ступенчатых в/з функций на $I$


есть просто множество всех кусочно-постоянных (к.-п.) и непрерывных


справа (н.~спр.) в/з функций на $I$. Целесообразно ввести


$\sigma$-алгебру ${\cal B}$ п/м $I$, порожденную п/а ${\cal I}$,


т.\,е. борелевское оснащение $I$. В этом и в следующем разделах


полагаем, что задана некоторая п/а


${\cal L}$ подмножеств (п/м) $I$,


для которой ${\cal I} \subset {\cal L} \subset


{\cal B}$. Промежуток $I$ используем здесь в качестве множества


$E$ (см. часть I). Постулируем, что все компоненты ${b_1} = {{b_1}(\cdot)},


\dotsc,{b_n} = {{b_n}(\cdot)}$ вектор-функции $b$ суть элементы


пространства $B(I,{\cal L})$ (см. раздел 5) равномерных пределов


${\cal L}$-ступенчатых в/з функций на $I$ (если ${\cal L} = {\cal I}$,


то ${b_1},\dotsc,{b_n}$ суть равномерные пределы некоторых


последовательностей к.-п. и н. спр. в/з функций на $I$; если


${\cal L} = {\cal B}$, то ${b_1},\dotsc,{b_n}$ --- ограниченные


борелевские функции на $I$). В качестве $f$ используем в


простейшем случае функции из ${B_0}(I,{\cal L})$


(см. раздел 5), хотя допускаем


и вариант применения функций из $B(I,{\cal L})$. Полагаем, что в (8.1)


задано (основное) начальное условие ${x(t_0)} = {x_0} \in


{{\Bbb R}^n}$, а


выбор $f$ стеснен некоторыми ограничениями, часть из которых


может определяться краевыми и промежуточными условиями, а также фазовыми


ограничениями (решение (8.1) понимаем в смысле Каратеодори [1]).


В этих условиях представляет интерес исследование пучка возможных


траекторий и его сечения гиперплоскостью $t = {{\vt}_0}$, т.\,е. области


достижимости (ОД) в момент ${{\vt}_0}$. Иногда возникает необходимость


в рассмотрении некоторых преобразований фазового вектора.


Oбсудим также такую возможность, полагая заданными натуральное число


$p \in {\cal N}$ и непрерывный оператор $\zeta$, действующий из


${\Bbb R}^n$ в ${{\Bbb R}^p}$. В этих условиях также можно говорить об


ОД (или о пучке траекторий), но только для значений


преобразования~$\zeta$.





Будем придерживаться содержательного способа изложения.


В этом и в следующем разделах


через ${\Phi}(\cdot,\cdot)$ обозначаем фундаментальную матрицу


решений (матрицант) ([1], [3], [26])


однородной системы $\Dot x(t) =


{A(t)}{x(t)}$, а через $\ld$ --- след меры Лебега--Бореля на ${\cal L}$.


Если ${t_*} \in {I_0}$ и ${x_*} \in {{\Bbb R}^n}$, то при


$f \in B(I,{\cal L})$ решение уравнения (8.1) (рассматриваемого на


$[{t_*},{{\vt}_0}]$) с начальным условием ${x(t_*)} = {x_*}$


определяется формулой Коши [1], [3], [26]


\begin{equation}


{x(t)} = {\Phi}(t,{t_*}){x_*} +


{\int_{[{t_*},t[}}{f(\tau)}{\Phi}(t,\tau){b(\tau)}{\ld}(d\tau)\quad


({t_*} \le t \le {{\vt}_0}),


\end{equation}


где интеграл вектор-функции определяется покомпонентно;


это решение обозначаем через


\begin{equation}


{{\vp}_f}(\cdot,{t_*},{x_*}) =


({{\vp}_f}(t,{t_*},{x_*}))_{t \in [{t_*},{{\vt}_0}]},


\end{equation}


получая непрерывную $n$-вектор-функцию на $[{t_*},{{\vt}_0}]$.


Среди свойств решения (8.1)--(8.3) отметим известное


полугрупповое свойство. Наконец, для $\zeta$-преобразования


(8.2), (8.3) естественно использовать соглашение: при


${t_*} \in {I_0}$, ${x_*} \in {{\Bbb R}^n}$, $f \in B(I,{\cal L})$


и $t \in [{t_*},{{\vt}_0}]$ значение этого $\zeta$-преобразованного


движения есть вектор ${\zeta}({{\vp}_f}(t,{t_*},{x_*})) \in {{\Bbb R}^p}$.


В частности, можно полагать ${t_*} = {t_0}$, ${x_*} = {x_0}$ и


$t = {{\vt}_0}$, получая точки ОД, если при этом соблюдается


соответствующая система ограничений.





Редакция формулы Коши, определяемая в (8.2), допускает естественное


расширение: в обозначениях первой части рассматриваем при


${t_*} \in {I_0}$, ${x_*} \in {{\Bbb R}^n}$ и $\mu \in {{\Bbb A}({\cal L})}$


функцию


\begin{equation}


{{\wt \vp}_\mu}(\cdot,{t_*},{x_*}) =


({{\wt \vp}_\mu}(t,{t_*},{x_*}))_{t \in [{t_*},{{\vt}_0}]}


\end{equation}


из $[{t_*},{{\vt}_0}]$ в ${{\Bbb R}^n}$, для которой имеет место


$\fo{t} \in [{t_*},{{\vt}_0}]$


\begin{equation}


{{\wt \vp}_\mu}(t,{t_*},{x_*}) \stac {\Phi}(t,{t_*}){x_*} +


{\int_{[{t_*},t[}}{\Phi}(t,\tau){b(\tau)}{{\mu}(d\tau)}.


\end{equation}





Можно рассматривать $n$-вектор-функцию (8.4), (8.5) как некое


обобщенное движение. Напомним, что (см. [1], [3], [26])


${\Phi}(\cdot,\cdot)$ обладает полугрупповым свойством. С учетом


этого и свойства конечной аддитивности неопределенного интеграла


легко проверяется





\begin{props}


Если ${t_*} \in {I_0}$, ${x_*} \in {{\Bbb R}^n}$, $\mu \in


{{\Bbb A}({\cal L})}$, ${t^*} \in [{t_*},{{\vt}_0}]$ и


$t \in [{t^*},{{\vt}_0}]$, то


$$


{{\wt \vp}_\mu}(t,{t_*},{x_*}) =


{{\wt \vp}_\mu}(t,{t^*},{{\wt \vp}_\mu}({t^*},{t_*},{x_*})).


$$


\end{props}





В самом деле, фиксируя ${t_*}$, ${x_*}$, $\mu$, ${t^*}$ и $t$


в соответствии с условиями доказываемого предложения, заметим


с учетом (8.5), что


\begin{multline*}


{{\wt \vp}_\mu}(t,{t^*},{{\wt \vp}_\mu}({t^*},{t_*},{x_*})) =


{\Phi}(t,{t^*}){{\wt \vp}_\mu}({t^*},{t_*},{x_*}) +


{\int_{[{t^*},t[}}{\Phi}(t,\tau){b(\tau)}{\mu}(d\tau) =\\


= {\Phi}(t,{t^*}){\Phi}({t^*},{t_*}){x_*} +


{\Phi}(t,{t^*}){\int_{[{t_*},{t^*}[}}{\Phi}({t^*},\tau){b(\tau)}{\mu}(d\tau)


+\\


+ {\int_{[{t^*},t[}}{\Phi}(t,\tau){b(\tau)}{\mu}(d\tau)


= {\Phi}(t,{t_*}){x_*} +


{\int_{[{t_*},{t^*}[}}{\Phi}(t,\tau){b(\tau)}{\mu}(d\tau) +\\


+ {\int_{[{t^*},t[}}{\Phi}(t,\tau){b(\tau)}{\mu}(d\tau) =


{\Phi}(t,{t_*}){x_*} +


{\int_{[{t_*},t[}}{\Phi}(t,\tau){b(\tau)}{\mu}(d\tau) =


{{\wt \vp}_\mu}(t,{t_*},{x_*}).


\end{multline*}





Итак, в (8.4), (8.5) определена динамическая система. Эти представления


дополняются преобразованием на основе $\zeta$, т.\,е. можно ввести обобщенные


траектории в виде зависимостей


\begin{equation}


t \longmapsto {\zeta}({{\wt \vp}_\mu}(t,{t_*},{x_*})):


[{t_*},{{\vt}_0}] \lo {{\Bbb R}^p},


\end{equation}


где ${t_*} \in {I_0}$, ${x_*} \in {{\Bbb R}^n}$ и $\mu \in


{{\Bbb A}({\cal L})}$. В этой связи следует учесть весьма очевидную


возможность погружения обычных траекторий, определяемых на основе


(8.2), в соответствующее пространство обобщенных траекторий. Ограничимся


случаем ${t_*} = {t_0}$ и ${x_*} = {x_0}$.


Следуем обозначениям [5], [6], [10] в части,


касающейся построения неопределенного


интеграла: $(I,{\cal L},\ld)$ есть вариант пространства (3.5.1) из [6] и


для $f \in B(I,{\cal L})$ можно рассматривать неопределенный


$\ld$-интеграл $f$ ([5], c.\,69), $f*\ld \in {{\Bbb A}({\cal L})}$.


Тогда (см. (8.2), (8.5) и свойства из ([5], c.\,69)) при $\mu = f*\ld$


\begin{equation}


{{\vp}_f}(\cdot,{t_*},{x_*}) = {{\wt \vp}_\mu}(\cdot,{t_*},{x_*}).


\end{equation}





\begin{notes}


Было рассмотрено погружение обычных управлений и траекторий в


соответствующее пространство обобщенных элементов (ОЭ),


следуя идеологии ([5], c.\,136), в основе которой --- преобразование


функции в меру посредством построения неопределенного $\ld$-интеграла.


Полезно, однако, отметить, что в (8.4), (8.5) допускается использование


в качестве управлений любых мер из ${\Bbb A}({\cal L})$, а это


позволяет обслуживать другие содержательные задачи. В частности,


можно отказаться от представления обычных траекторий посредством


(8.2), (8.3) и обратиться к ситуации управления с толчками [12],


используя уже схему построения ОЭ, рассматриваемую в ([6], гл.\,7).


\end{notes}





Возвращаясь к (8.6), отметим, что $\zeta$-траектории зависят от


управления $\mu \in {{\Bbb A}({\cal L})}$ непрерывно в смысле


$*$-слабой топологии ${\Bbb A}({\cal L})$ и топологии поточечной


сходимости пространства всех $n$-вектор-функций


[13]. Это свойство вытекает из (8.5). Oграничимся, однако,


замечанием об отображении


\begin{equation}


\mu \longmapsto {\zeta}({{\wt \vp}_\mu}({{\vt}_0},{t_0},{x_0})):


{{\Bbb A}({\cal L})} \lo {{\Bbb R}^p},


\end{equation}


имея в виду проблему корректного расширения задачи о построении


ОД; разумеется, ${\Bbb A}({\cal L})$ оснащаем при этом топологией


${{\tau}_*}({\cal L})$, а ${\Bbb R}^p$ --- обычной топологией


${{\tau}_{\Bbb R}^{(p)}}$ покоординатной сходимости.


Имеем в (8.8) непрерывное в указанном смысле отображение. В свою


очередь, зависимость обобщенной траектории от $\mu \in {{\Bbb A}({\cal L})}$


также непрерывна. Здесь целесообразно провести более подробное обсуждение


топологических конструкций, имея в виду применения в схеме раздела 5


при ${\cal R} \stac {{\Bbb R}^n}$. Пространство ${\cal R}^{I_0}$ всех


$n$-вектор-функций на $I_0$ оснащаем топологией


${{\otimes}^{I_0}}({\tau}_{\Bbb R}^{(n)})$ тихоновского произведения


экземпляров $({\cal R},{{\tau}_{\Bbb R}^{(n)}})$ с индексным множеством


${I_0}$, т.\,е. обычной топологией поточечной сходимости. Тогда


\begin{equation}


\mu \longmapsto {{\wt \vp}_\mu}(\cdot,{t_0},{x_0}):


{{\Bbb A}({\cal L})} \lo {{\cal R}^{I_0}}


\end{equation}


есть непрерывный в смысле топологий


\begin{equation}


{{\tau}_*}({\cal L}),\ \ {{\otimes}^{I_0}}({\tau}_{\Bbb R}^{(n)})


\end{equation}


оператор [13], что непосредственно следует из (8.5) и определения


${{\tau}_*}({\cal L})$.


Итак, (8.9) есть непрерывный в смысле ТП (8.10) оператор. Вернемся


к (8.2), (8.3), (8.7). Выбор управления $f \in B(I,{\cal L})$


стеснен обычно теми или иными ограничениями. Один из наиболее естественных


случаев связан с ресурсным ограничением


\begin{equation}


{\int_I}|{f(t)}|{{\ld}(dt)} \le c,


\end{equation}


где $c \in [0,\infty[$. Разумеется, возможны и другие ограничения, связанные,


в частности, с осуществлением краевых и промежуточных условий. Сейчас, однако,


ограничимся обсуждением (8.11). В этом случае, ориентируясь на (8.7),


отметим, что множество $F_c$ всех функций $f \in B(I,{\cal L})$


со свойством (8.11) в виде всюду плотного подмножества (п/м)


погружается в некоторое множество


$$


{K_c} \in ({{\tau}_*}({\cal L}) - {\comp})[{\Bbb A}({\cal L})] \setminus


\{\emp\},


$$


которое определяется в ([5], cc.\,82, 87; [6], c.\,54). Следуя ([6], c.\,39),


введем для каждой меры $\mu \in {{\Bbb A}({\cal L})}$ вариацию


${v_\mu} \in {(\add)_+}[{\cal L}]$ как функцию множеств, а также


множество ${{\Bbb A}_\ld}[{\cal L}]$ всех $\mu \in {{\Bbb A}({\cal L})}$


со свойством


$\fo{L} \in {\cal L}$


\begin{equation}


({{\ld}(L)} = 0) \Lo ({{\mu}(L)} = 0).


\end{equation}


Тогда полагаем ${K_c} \stac \{\mu \in {{\Bbb A}_\ld}[{\cal L}] \mid


{v_\mu}(I) \le c\}$. В этом случае для топологий


${{\tau}_*}({\cal L})$, ${{\tau}_0}({\cal L})$, ${{\tau}_\otimes}({\cal


L})$ и ${{\tau}_{\bold B}^*}({\cal L})$,


определяемых в ([6], гл.\,3) и составляющих


множество ${{\frak M}({\cal L})} \stac \{{{\tau}_*}({\cal L});


{{\tau}_0}({\cal L});{{\tau}_\otimes}({\cal L});


{{\tau}_{\bold B}^*}({\cal L})\}$,


имеем свойство


$\fo{\tau} \in {{\frak M}({\cal L})}$


\begin{equation}


{K_c} = {\cl}(\{f*\ld:f \in {F_c}\},\tau).


\end{equation}





Как следствие получим для пучка обычных траекторий плотное погружение


в аналогичный пучок обобщенных траекторий


\begin{equation}


\{{{\wt \vp}_\mu}(\cdot,{t_0},{x_0}):\mu \in {K_c}\} =


{\cl}(\{{{\vp}_f}(\cdot,{t_0},{x_0}):f \in {F_c}\},


{{\otimes}^{I_0}}({\tau}_{\Bbb R}^{(n)})).


\end{equation}





Условие (8.12) именуют слабой абсолютной непрерывностью относительно


$\ld$; (8.13) --- типичное для [5], [6], [10] расширение пространства


управлений (универсальное в диапазоне топологий); (8.14)


подобно в логическом отношении свойствам, рассматриваемым в


([1], гл.\,III, IV) для нелинейных задач управления с геометрическими


ограничениями. Сейчас ограничимся достаточно простым случаем


промежуточных условий, дополняющих требования на выбор


$f \in {F_c}$. Именно, пусть $Y$ --- замкнутое в $({\cal R},


{{\tau}_{\Bbb R}^{(n)}})$ подмножество ${\cal R}$ и ${t_1} \in {I_0}$;


требуется обеспечить свойство ${{\vp}_f}({t_1},{t_0},{x_0}) \in Y$.


При этом $Y$-ограничении (и при условии (8.11)) требуется построить


регуляризированный вариант ОД. Уже в этом достаточно простом случае


имеем нелинейную задачу раздела 2 при следующих соглашениях


относительно параметров: ${\bold F} = {F_c}$, ${\bold X} = {\cal R} =


{{\Bbb R}^n}$, ${\bold H} = {{\Bbb R}^p}$, ${{\bold s}(f)} =


{{\vp}_f}({t_1},{t_0},{x_0})$ и ${{\bold h}(f)} =


{\zeta}({{\vp}_f}({\vt}_0,{t_0},{x_0}))$ при


$f \in {\bold F}$, ${\bold Y} = Y$.


В части описания модели раздела 4 заметим, что множество ${\bold K}$


отождествляется здесь с ${K_c}$, что соответствует соглашениям разделов


5 и 6. Мы соблюдаем соглашения о выборе топологий $t_{\bold l}$ и


$t_{\bold u}$ множества ${K_c}$. Однако вопрос о конкретизации


$\mG$ и о соответствующей этому выбору конкретизации ${\bold X}$


сейчас в общей постановке не рассматриваем, поскольку важная


в этом смысле детализация ${\mG}_0$ (см. (5.8))


обычно определяется конкретными условиями решаемой задачи.


Не вникая в логику построения битопологических пространств первой части,


отметим, что объективно отображения $g$ и $\omega$ раздела


4 можно определить соответственно в виде зависимостей


$$


\mu \longmapsto {{\wt \vp}_\mu}({t_1},{t_0},{x_0}),\ \


\mu \longmapsto {\zeta}({{\wt \vp}_\mu}({{\vt}_0},{t_0},{x_0})),


$$


определенных на ${K_c}$. Тогда множество ${{\omega}^1}({g^{-1}}({\bold Y}))$


приобретает смысл ОД обобщенной управляемой системы в условиях


ограничений на выбор управления $\mu \in {{\Bbb A}_\ld}[{\cal L}]$,


определяемого: 1)~$c$-ограничением на полную вариацию,


2)~${\bold Y}$-ограничением на реализацию вектора


${{\wt \vp}_\mu}({t_1},{t_0},{x_0})$. Более тонкий анализ,


связанный с представлением $g$ и построением битопологического


пространства $({\bold X},{{\tau}_{\bold l}},{{\tau}_{\bold u}})$,


проведем в следующем разделе для простейшего случая задачи управления


материальной точкой.





\section{Одна задача управления материальной точкой на плоскости}





В этом разделе для частного вида системы (8.1)


\begin{equation}


{{\Dot x}_1}(t) = {x_3}(t),\ \ {{\Dot x}_2}(t) =


{x_4}(t),\ \ {{\Dot x}_3}(t) = {{b_1}(t)}{f(t)},\ \


{{\Dot x}_4}(t) = {{b_2}(t)}{f(t)}


\end{equation}


рассматриваем, как и ранее, процесс управления на $I_0$


с заданным н.\,у. ${x(t_0)} = {x_0} \in {{\Bbb R}^4}$ (можно было


бы ограничиться случаем ${t_0} = 0$ и ${{\vt}_0} = 1$);


полагаем, что ${b_1}(\cdot) \in {B_0}(I,{\cal L})$ и


${b_2}(\cdot) \in {B_0}(I,{\cal L})$. Для простоты полагаем, что в


случае $n = 4$ множество $Y$ имеет смысл произведения


замкнутых множеств ${Y_1}$, ${Y_1} \subset {\Bbb R} \times {\Bbb R}$,


и ${Y_2}$, ${Y_2} \subset {\Bbb R} \times {\Bbb R}$.


Множество $Y_1$ характеризует условие на реализацию


$({{x_1}(t_1)},{{x_2}(t_1)})$, а множество $Y_2$ --- условие на реализацию


$({{x_3}(t_1)},{{x_4}(t_1)})$:


\begin{equation}


(({{x_1}(t_1)},{{x_2}(t_1)}) \in {Y_1}) \& (({{x_3}(t_1)},{{x_4}(t_1)}) \in


{Y_2}).


\end{equation}





Разумеется, в качестве $x(\cdot)$ в (9.2) следует использовать


траекторию, определяемую в (9.1), т.\,е. вектор-функцию вида (8.3).


Используя известную конкретизацию ${\Phi}(\cdot,\cdot)$


[3], [26], получаем, что (9.2) преобразуется к виду


\begin{multline*}


\bigg(\bigg({\int_{[{t_0},{t_1}[}}({t_1} - t){{b_1}(t)}{f(t)}{\ld}(dt),


{\int_{[{t_0},{t_1}[}}({t_1} - t){{b_2}(t)}{f(t)}{\ld}(dt)\bigg) \in


{{\wt Y}_1}\bigg) \&\\


\& \bigg(\bigg({\int_{[{t_0},{t_1}[}}{{b_1}(t)}{f(t)}{\ld}(dt),


{\int_{[{t_0},{t_1}[}}{{b_2}(t)}{f(t)}{\ld}(dt)\bigg) \in {{\wt Y}_2}\bigg),


\end{multline*}


где множества ${{\wt Y}_1}$ и ${{\wt Y}_2}$


получаются сдвигом $Y_1$ и ${Y_2}$ соответственно; однако для простоты


полагаем далее, что все компоненты $x_0$ нулевые, а тогда


${{\wt Y}_1} = {Y_1}$ и ${{\wt Y}_2} = {Y_2}$;


пусть, кроме того, ${t_0} = 0$, т.\,е. в нашем случае


${I_0} = [0,{{\vt}_0}]$.


Мы не конкретизируем отображение $\zeta$ раздела 8.


Здесь несколько изменим параметры (в сравнении с разделом 8), полагая


далее в этом разделе $n = 2$ (имеется в виду конкретизация размерности


в разделе 5); пусть $\mG = \{1;2\}$ (неупорядоченная пара) и


${{\mG}_0} = \{2\}$ (одноэлементное множество). Тогда


${\cal R} = {{\Bbb R}^2}$. Матрицант (5.7) конкретизируем следующим


образом, полагая, что ${{\chi}_{[{t_0},{t_1}[}} \in {{\Bbb R}^I}$


есть индикатор промежутка $[{t_0},{t_1}[$; именно, при $t \in I$


полагаем


\begin{multline}


({S_{1,1}}(t) = ({t_1} - t){b_1}(t){{\chi}_{[{t_0},{t_1}[}}(t)) \&


({S_{2,1}}(t) = ({t_1} - t){b_2}(t){{\chi}_{[{t_0},{t_1}[}}(t)) \&\\


\& ({S_{1,2}}(t) = {b_1}(t){{\chi}_{[{t_0},{t_1}[}}(t)) \&


({S_{2,2}}(t) = {b_2}(t){{\chi}_{[{t_0},{t_1}[}}(t)).


\end{multline}


Далее, отображение $g$ определяется в (5.11) при конкретизации


(9.3). Напомним, что ${\bold X} = {{\cal R}^{\mG}}$,


а ${\bold Y}$ есть п/м ${\bold X}$, для которого


$$


{\bold Y} = \{y \in {{\cal R}^{\mG}} \mid ({y(1)} \in {Y_1}) \&


({y(2)} \in {Y_2})\}.


$$


В этих условиях требование (9.2) сводится к ${\bold Y}$-ограничению части I,


реализуемому в терминах (5.11). Погружение ${\bold F} = {F_c}$ в


${\bold K} = {K_c}$ осуществляем на основе построения неопределенного


$\ld$-интеграла $f \in {\bold F}$.


Полагаем ${S_{1,1}} \notin {B_0}(I,{\cal L})$ и


${S_{2,1}} \notin {B_0}(I,{\cal L})$, что непременно имеет


место при ${t_1} > 0$ и естественных условиях


невырожденности ${b_1}(\cdot)$ и ${b_2}(\cdot)$;


например, можно полагать, что последние две функции к.-п., н.~спр.


и не обращаются в нуль на $[{t_0},{t_1}[$.


Тогда (в силу (9.3))


определение ${\mG}_0$ соответствует (5.8). Отождествляем


${\bold H}$ с ${{\Bbb R}^p}$, где $p \in {\cal N}$, а


${\bold h}$ определяем как


\begin{equation}


f \longmapsto {\zeta}({{\vp}_f}({{\vt}_0},0,{\Bbb O})):{F_c} \lo


{{\Bbb R}^p},


\end{equation}


где ${\Bbb O} \in {{\Bbb R}^4}$ --- вектор, все координаты которого


нулевые. Разумеется, $\zeta$ соответствует разделу 8 и переводит


${\Bbb R}^4$ в ${{\Bbb R}^p}$ (непрерывность $\zeta$ постулируется).


Следуя конструкциям части I, можно ввести несколько вариантов ослабления


${\bold Y}$-ограничения (в силу конечности $\mG$ определения


упрощаются). Oграничимся


содержательным обсуждением; восстановление необходимых


для строгого изложения деталей не составляет труда и


фактически повторяет [5], [6]. Итак, в первом случае, соответствующем


в существенной части определению (5.6), вместо $Y_1$ и $Y_2$


(в (9.2)) используем их $\var$-окрестности. Во втором случае, следуя


(5.17) и (5.18), вводим $\var$-окрестность только для


${Y_1}$. Впрочем, можно было бы ввести топологию


$\tau_{\bold u}$ (5.10), расширяя спектр возможных вариантов


ослабления условия (9.2). Во всех этих случаях ослабленные версии


содержательного условия (9.2) трансформируются в соответствующие


ослабленные аналоги ${\bold Y}$-ограничения.


Соответствующие варианты ${{\cal Y}_{\bold l}}$, ${{\cal Y}_{\bold u}}$


и ${\wt {\cal Y}}_{\bold u}$ обслуживают, таким образом, весьма


содержательные и важные, с практической точки зрения, ослабленные аналоги


(9.2); в результате реализуются конкретные варианты (5.16), предложения 5.2


и теоремы 6.1. Ограничимся определением


МП (5.16). В этой связи полезно конкретизировать


(8.5) для случая (9.1): если $\mu \in {{\Bbb A}({\cal L})}$, то траектория


(8.4), (8.5) при нулевых начальных условиях характеризуется четырьмя


функциями, для которых при $t \in {I_0}$


\begin{alignat*}{2}


{{\wt \vp}_{\mu,1}}(t,0,{\Bbb O}) &=


{\int_{[0,t[}}(t - \tau){b_1}(\tau){\mu}(d\tau), &\quad


{{\wt \vp}_{\mu,2}}(t,0,{\Bbb O}) &=


{\int_{[0,t[}}(t - \tau){b_2}(\tau){\mu}(d\tau),\\


\qquad {}{{\wt \vp}_{\mu,3}}(t,0,{\Bbb O}) &=


{\int_{[0,t[}}{b_1}(\tau){\mu}(d\tau), &\quad


{{\wt \vp}_{\mu,4}}(t,0,{\Bbb O}) &= {\int_{[0,t[}}{b_2}(\tau){\mu}(d\tau).


\end{alignat*}


В этих соотношениях наиболее важны случаи $t = {t_1}$ и


$t = {{\vt}_0}$. В терминах двух вышеупомянутых моментов


времени вводим задачу об определении множества


${\Bbb G}_\partial$ всех векторов ${\zeta}({{\wt \vp}_\mu}({{\vt}_0},


0,{\Bbb O}))$ таких, что $\mu \in {K_c}$ и при этом


\begin{equation}


(({{\wt \vp}_{\mu,1}}({t_1},0,{\Bbb O}),{{\wt \vp}_{\mu,2}}({t_1},


0,{\Bbb O})) \in {Y_1}) \& (({{\wt \vp}_{\mu,3}}({t_1},0,{\Bbb O}),


{{\wt \vp}_{\mu,4}}({t_1},0,{\Bbb O})) \in {Y_2}).


\end{equation}





В (9.5) дано естественное расширение (9.2). В терминах (9.5)


получили задачу определения ОД обобщенной системы; эту задачу


можно истолковать в терминах условия ${g(\mu)} \in {\bold Y}$


(см. раздел 5). Достаточно учесть (9.3) и определение


(5.11). Разумеется, (9.5) реализует в виде множества


допустимых элементов соответствующий конкретный вариант


${g^{-1}}({\bold Y})$. Что же касается $\omega$,


то данное отображение есть расширенная версия (9.4), т.е.


\begin{equation}


\mu \longmapsto {\zeta}({{\wt \vp}_\mu}({{\vt}_0},0,{\Bbb O})):


{\bold K} \lo {{\Bbb R}^p};


\end{equation}


связь (9.4), (9.6)


реализуется в терминах ${\bold m} \in {\bold K}^{\bold F}$,


${\bold m}(f) = {f*\ld}$ (см. [6], теоремa 3.7.5). В виде


${\Bbb G}_\partial$ имеем теперь требуемую версию


${{\omega}^1}({g^{-1}}({\bold Y}))$, т.\,е. универсальное в смысле раздела


5 МП (см. (5.16)); при этом из определения ${\mG}_0$


видно, что наша асимптотическая версия задачи  о построении ОД


обладает определенной грубостью при ослаблении условий на реализацию


вектора скорости в (9.2) (имеется в виду ограничение,


определяемое в терминах $Y_2$). Указанное свойство механического


характера подобно отмечавшимся в [5], [6], [8], [10] свойствам грубости при


ослаблении части ограничений. Рассматриваемую здесь постановку,


связанную с корректным расширением задачи о построении ОД


для системы (9.1), можно существенно усложнить, сохраняя в той или


иной форме упомянутое свойство грубости в отношении соблюдения


условий на реализацию вектора скорости как функции времени.





\section{О некоторых свойствах плотности}





Возвратимся к рассмотрению общих конструкций первой части.


Следуем соглашениям раздела 5 в части обозначений, связанных с


измеримым пространством (ИП) $(E,{\cal L})$, где $E$ ---


непустое множество, а ${\cal L}$ --- п/а п/м $E$. Итак,


$(E,{\cal L})$ есть ИП с п/а множеств. В разделах 5 и 7


показано, что ${\frak K}$ (5.2) весьма представительно


как семейство п/м ${\Bbb A}({\cal L})$.


Этим мотивировано требование ${\bold K} \in {\frak K}$


раздела 5. Сейчас обсудим некоторые аспекты, связанные


с условием 5.1; уже отмечались (см. замечание 5.1)


различные частные случаи общей постановки раздела 5, для которых это


условие выполняется [5], [6], [10].


Pассмотрим еще один гипотетический случай, когда предположение


о плотности, постулируемое в условии 5.1, имеет место.


Однако это суждение будет дано в более общей форме --- в виде


свойства плотности, универсального в смысле топологий


${{\Bbb A}({\cal L})}$, используемых в ([6], гл.\,3; [10]).


Условимся о традиционном обозначении, понимаемом


здесь более широко, чем в разделе 7: если $X$ --- непустое множество,


$A$ и $B$ --- суть п/м ${{\Bbb R}^X}$, то


$$


A \oplus B \stac \{a + b:(a,b) \in A \times B\} \in


{\cal P}({\Bbb R}^X)


$$


есть алгебраическая сумма $A$, $B$ [1], [11]. В качестве


$X$ используем обычно $E$ или ${\cal L}$. Известно [5], [6], [10],


что применение


к.-а. мер в качестве обобщенных элементов (ОЭ)


по схеме [5], [6] весьма плодотворно в следующих двух


случаях: 1)~в условиях, когда пространство обычных решений


предкомпактно в смысле $*$-слабой топологии ${\Bbb A}({\cal L})$


(при надлежащем погружении в пространство упомянутых ОЭ);


2)~в случае, когда обычные решения --- в/з функции на $E$ ---


неотрицательны.


Поэтому представляется естественным рассмотрение некой гибридной,


в смысле 1),~2), конструкции. Используем здесь простейшие свойства


топологических векторных пространств (ТВП) и положения [5], [6].


Итак, рассматриваем в качестве обычных решений функции на $E$


и исследуем свойства плотности в духе условия 5.1. Более того,


рассматриваем и другие свойства плотности, подобные тем, что использовались


в [8], [10] при исследовании МП в пространстве


ОЭ: речь идет об аналогах предложений 4.1, 4.2. Изучение таких


промежуточных МП на предмет универсальности представлений в терминах


множества ${g^{-1}}({\bold Y})$ проясняет структуру


асимптотической версии задачи о достижимости (здесь имеется в виду


универсальность в смысле топологического оснащения пространства


ОЭ). В этой связи рассматриваем вопрос об аналогах условия 5.1


более широко.





Через ${B_0^+}(E,{\cal L})$ и ${B^+}(E,{\cal L})$ обозначим конусы


всех неотрицательных (поточечно) элементов линейных пространств


${B_0}(E,{\cal L})$ и $B(E,{\cal L})$ соответственно (см. раздел 5).


Если $f \in {{\Bbb R}^E}$, то через $\bmid{f}\bmid$ обозначим


в/з функцию на $E$ со значениями $|{f(x)}|$, $x \in E$. При этом


$\bmid{f}\bmid \in {B_0^+}(E,{\cal L})$\, $\fo{f} \in {B_0}(E,{\cal L})$.


Аналогичным образом, $\bmid{g}\bmid \in {B^+}(E,{\cal L})$\,


$\fo{g} \in B(E,{\cal L})$. Фиксируем $\eta \in {(\add)_+}[{\cal L}]$ и


$\kappa \in [0,\infty[$;


введем также


\begin{equation}


\bigg(U \stac \bigg\{f \in {B_0}(E,{\cal L})\Big|


{\int_E}\bmid{f}\bmid{d\eta} \le


\kappa\bigg\}\bigg) \& \bigg(V \stac \bigg\{f \in B(E,{\cal L}) \Big|


{\int_E}\bmid{f}\bmid{d\eta} \le \kappa\bigg\}\bigg).


\end{equation}





Множества (10.1) подобны множеству $F_c$ раздела 8 и соответствуют


естественным в теории управления ресурсным ограничениям (см. (8.11)).


Рассмотрим следующие два варианта множества ${\bold F}$:


$$


{\bold F} = {B_0^+}(E,{\cal L}) \oplus U,\ \ {\bold F} =


{B^+}(E,{\cal L}) \oplus V.


$$


В первом случае получаем непустое п/м ${B_0}(E,{\cal L})$, а во


втором --- непустое п/м $B(E,{\cal L})$. Элементы этих множеств


рассматриваем в качестве обычных управлений. Потребуется


также одно специальное множество в пространстве ${{\Bbb A}({\cal L})}$.


Пусть [5], [6], [10] ${(\add)^+}[{\cal L};\eta]$ --- множество всех


$\mu \in {(\add)_+}[{\cal L}]$ таких, что


$\fo{L} \in {\cal L}$


\begin{equation}


({{\eta}(L)} = 0) \Lo ({{\mu}(L)} = 0);


\end{equation}


в (10.2) используется свойство слабой абсолютной $\eta$-непрерывности


[27] ((10.2) --- абстрактный аналог условия (8.12)).


При этом ${(\add)^+}[{\cal L};\eta]$ --- конус в


${{\Bbb A}({\cal L})}$; введем порожденное этим конусом


линейное подпространство (п/п) ${{\Bbb A}({\cal L})}$, обозначаемое


через ${{\Bbb A}_\eta}[{\cal L}]$ ([6], c.\,43). Далее, следуя ([6], c.\,40),


введем сильную норму ${{\Bbb A}({\cal L})}$, определяемую


посредством полной вариации: если $\mu \in {{\Bbb A}({\cal L})}$, то


([6], cc.\,39, 40) ${{v_\mu}(E)} \in [0,\infty[$ есть def полная вариация


$\mu$ на множестве $E$. Введем


\begin{equation}


W \stac \{\mu \in {{\Bbb A}_\eta}[{\cal L}] \mid {v_\mu}(E) \le \kappa\}.


\end{equation}


Разумеется, $W$ есть слабо абсолютно непрерывная (относительно


$\eta$) часть шара в ${{\Bbb A}({\cal L})}$. Возьмем непустое


множество


\begin{equation}


{(\add)^+}[{\cal L};\eta] \oplus W


\end{equation}


в качестве ${\bold K}$. Если $f \in B(E,{\cal L})$, то $f*\eta \in


{{\Bbb A}_\eta}[{\cal L}]$ есть def неопределенный $\eta$-интеграл


$f$ ([5], cc.\,69, 86). Как и в [5], [6], [10], [13], рассматриваем погружение


$B(E,{\cal L})$ в ${{\Bbb A}_\eta}[{\cal L}]$ посредством оператора


\begin{equation}


f \longmapsto {f*\eta}:B(E,{\cal L}) \lo {{\Bbb A}_\eta}[{\cal L}].


\end{equation}


Один из вариантов такого погружения использовался в разделе 8.


В частности, можем рассматривать свойства образов множеств


$$


{B_0^+}(E,{\cal L}) \oplus U,\ \ {B^+}(E,{\cal L}) \oplus V


$$


при действии оператора (10.5). Отметим ряд свойств множества (10.4),


связанных с конструкциями раздела 7. В силу известных положений


теории топологических групп ([11], c.\,449)


множество (10.4) замкнуто в локально выпуклом ТВП


$({{\Bbb A}({\cal L})},{{\tau}_*}({\cal L}))$;


здесь используется свойства $*$-слабой замкнутости


${(\add)^+}[{\cal L};\eta]$ и $*$-слабой компактности $W$ ([6], гл.\,3).


Кроме того,


$$


{(\add)^+}[{\cal L};\eta] \oplus W \in {{\frak C}'_*}[{\cal


L}],


$$


что непосредственно следует из положений раздела 7.


С учетом свойства транзитивности операции перехода к п/п, а также


соотношений для ${{\tau}_\otimes}({\cal L})$ и ${{\tau}_*}({\cal L})$


в ([5], c.\,80; [6], c.\,45), имеем в обозначениях


разделов 3 и 7 вложение


\begin{equation}


({{\tau}_*}({\cal L}) - {\comp})[{\Bbb A}({\cal L})] \subset


({{\otimes}^{\cal L}}({{\tau}_{\Bbb R}}) - {\comp})[{{\Bbb R}^{\cal L}}].


\end{equation}





Из (10.3) и (10.6) следует, в частности, что $W \in


({{\otimes}^{\cal L}}({\tau}_{\Bbb R}) - {\comp})[{\Bbb R}^{\cal L}]$,


а поэтому (снова с использованием свойств топологических групп в


([11], c.\,449)) получаем:


множество ${(\add)^+}[{\cal L};\eta] \oplus W$ замкнуто в


$({{\Bbb R}^{\cal L}},{{\otimes}^{\cal L}}({\tau}_{\Bbb R}))$.





Ниже используются топологии ${{\tau}_*}({\cal L})$,


${{\tau}_\otimes}({\cal L})$,


${{\tau}_0}({\cal L})$ и ${{\tau}_{\bold B}^*}({\cal L})$ множества


${\Bbb A}({\cal L})$, определенные в ([6], гл.\,3); полагаем,


как и в частном случае ИП раздела 8, что


$$


{\frak M}({\cal L}) \stac \{{{\tau}_*}({\cal L});\;


{{\tau}_\otimes}({\cal L});\;{{\tau}_0}({\cal L});\;


{{\tau}_{\bold B}^*}({\cal L})\}.


$$





Наряду с ${{\otimes}^{\cal L}}({\tau}_{\Bbb R})$ оснащаем множество


${\Bbb R}^{\cal L}$ топологией ${{\otimes}^{\cal L}}({\tau}_\partial)$


тихоновского произведения экземпляров ${\Bbb R}$ с


(дискретной) топологией ${{\tau}_\partial} = {{\cal P}({\Bbb R})}$ каждого


экземпляра; в качестве


индексного множества здесь используется ${\cal L}$.


В силу ранее упомянутых свойств замкнутости


в топологиях ${{\tau}_*}({\cal L})$ и


${{\otimes}^{\cal L}}({\tau}_{\Bbb R})$ имеем, что


множество ${(\add)^+}[{\cal L};\eta] \oplus W$ замкнуто в


$({{\Bbb R}^{\cal L}},{{\otimes}^{\cal L}}({\tau}_\partial))$ и в


$({{\Bbb A}({\cal L})},{{\tau}_{\bold B}^*}({\cal L}))$.


Здесь достаточно учесть хорошо


известные свойства ([6], с.\,45; [11], с.\,463)


$$


({{\otimes}^{\cal L}}({\tau}_{\Bbb R}) \subset


{{\otimes}^{\cal L}}({\tau}_\partial)) \&


({{\tau}_*}({\cal L}) \subset {{\tau}_{\bold B}^*}({\cal L})).


$$





\begin{props}


Множество ${(\add)^+}[{\cal L};\eta] \oplus W$ замкнуто в ТП


$({{\Bbb A}({\cal L})},{{\tau}_0}({\cal L}))$


и \linebreak $({{\Bbb A}({\cal L})},{{\tau}_\otimes}({\cal L}))$.


\end{props}





Доказательство следует из уже упоминавшихся соотношений для топологий из


${\frak M}({\cal L})$ ([6], c.\,45); в частности, используется тот


факт, что ${{\tau}_\otimes}({\cal L})$ слабее каждой из


топологий ${{\tau}_*}({\cal L})$, ${{\tau}_0}({\cal L})$. Следует учесть,


кроме того, отмечавшееся ранее свойство замкнутости множества (10.4) в смысле


${{\otimes}^{\cal L}}({\tau_{\Bbb R}})$. Заметим, что особое значение для нас


имеет замкнутость множества (10.4) в топологии ${{\tau}_0}({\cal L})$,


что связано с условием 5.1.





Отметим, что (10.4) --- непустое п/м


${{\Bbb A}_\eta}[{\cal L}]$. Вернемся к вышеупомянутым версиям


${\bold F}$;


$$


{B_0^+}(E,{\cal L}) \oplus U \subset {B^+}(E,{\cal L}) \oplus V.


$$


Очевидно вложение $\{f*\eta:f \in {B^+}(E,{\cal L}) \oplus V\} \subset


{(\add)^+}[{\cal L};\eta] \oplus W$. Как следствие,


$\{f*\eta:f \in {B_0^+}(E,{\cal L}) \oplus U\} \subset


{(\add)^+}[{\cal L};\eta] \oplus W$.





\begin{props}


Имеет место


${(\add)^+}[{\cal L};\eta] \oplus W = {\cl}(\{{f*\eta}:f \in


{B_0^+}(E,{\cal L}) \oplus U\},{{\tau}_0}({\cal L})) =


{\cl}(\{{f*\eta}:f \in {B^+}(E,{\cal L}) \oplus V\},


{{\tau}_0}({\cal L}))$.


\end{props}





\begin{pf}


С учетом предложения 10.1 имеем,


что достаточно установить вложение


\begin{equation}


{(\add)^+}[{\cal L};\eta] \oplus W \subset {\cl}(\{{f*\eta}:


f \in {B_0^+}(E,{\cal L}) \oplus U\},{{\tau}_0}({\cal L})).


\end{equation}


Пусть (в этом доказательстве) $\ld$ --- элемент


множества в левой части (10.7), а


${{\ld}_1} \in {(\add)^+}[{\cal L};\eta]$ и ${{\ld}_2} \in W$


реализуют представление ${\ld} = {{\ld}_1} + {{\ld}_2}$.


С учетом положений ([6], гл.\,3) имеем


\begin{equation}


({{\ld}_1} \in {\cl}(\{{f*\eta}:f \in {B_0^+}(E,{\cal L})\},


{{\tau}_0}({\cal L}))) \&


({{\ld}_2} \in {\cl}(\{{f*\eta}:f \in U\},{{\tau}_0}({\cal L}))).


\end{equation}


Используем свойства базиса


$({\Bbb A}({\cal L}),{{\tau}_0}({\cal L}))$ ([5], c.\,81; [6], c.\,45).


Полагаем, что ${\Fin}({\cal L})$ есть def семейство всех непустых конечных


п/м ${\cal L}$. Для $\mu \in {{\Bbb A}({\cal L})}$ и ${\cal K} \in


{\Fin}({\cal L})$ рассматриваем множество ${T_\partial}(\mu,{\cal K})$ всех


$\nu \in {{\Bbb A}({\cal L})}$ таких, что $\fo{L} \in


{\cal K}$


$$


{{\mu}(L)} = {{\nu}(L)}.


$$


По свойствам ${{\tau}_0}({\cal L})$


[5], [6] имеем при $H \in {\cal P}({\Bbb A}({\cal L}))$ равенство


\begin{equation}


{\cl}(H,{{\tau}_0}({\cal L})) = \{\mu \in {{\Bbb A}({\cal L})} \mid


{T_\partial}(\mu,{\cal K}) \cap


H \ne \emp


\ \fo{\cal K} \in {\Fin}({\cal L})\}.


\end{equation}


В качестве $H$ можно использовать ${B_0^+}(E,{\cal L}) \oplus U$,


${B_0^+}(E,{\cal L})$ и $U$, т.\,е. комбинацию (10.8), (10.9).


Фиксируем ${\Bbb P} \in {\Fin}({\cal L})$. Тогда из (10.8), (10.9) имеем


\begin{equation}


({T_\partial}({{\ld}_1},{\Bbb P}) \cap \{{f*\eta}:f \in


{B_0^+}(E,{\cal L})\} \ne \emp) \& ({T_\partial}({{\ld}_2},{\Bbb P}) \cap


\{{f*\eta}:f \in U\} \ne \emp).


\end{equation}





С учетом (10.10) подбираем $\vp \in {B_0^+}(E,{\cal L})$ со свойством


${{\vp}*\eta} \in {T_\partial}({{\ld}_1},{\Bbb P})$. Кроме того, выбираем


(используя (10.10)) $\psi \in U$ со свойством ${{\psi}*\eta} \in


{T_\partial}({{\ld}_2},{\Bbb P})$. Тогда $\vp + \psi \in {B_0^+}(E,


{\cal L}) \oplus U$ реализует к.-а. меру


$$


(\vp + \psi)*\eta = ({\vp}*\eta) + ({\psi}*\eta) \in


{T_\partial}(\ld,{\Bbb P}).


$$


Поскольку выбор ${\Bbb P}$ был произвольным, установлено свойство:


$\fo{\cal K} \in {\Fin}({\cal L})$


\begin{equation}


{T_\partial}(\ld,{\cal K}) \cap


\{{f*\eta}:f \in {B_0^+}(E,{\cal L}) \oplus U\} \ne \emp.


\end{equation}





Из (10.9), (10.11) имеем


$\ld \in {\cl}(\{{f*\eta}:f \in {B_0^+}(E,{\cal L})


\oplus U\},{{\tau}_0}({\cal L}))$, чем и завершается обоснование


вложения (10.7).


\end{pf}





\begin{notes}


Из предложения 10.2 имеем вариант условия 5.1 (см. (5.5)),


где ${\bold F}$ конкретизировано


выше в двух вариантах, ${\bold K}$ отождествляется с множеством (10.4),


а отображение ${\bold m}$ раздела 5 реализуется в виде надлежащего


сужения (10.5) (сужение на ${\bold F}$). Итак, цель построения новой


конкретной версии конструкции раздела 5 фактически достигнута;


однако имеет смысл отметить ряд дополнительных свойств, имея


в виду более специализированные, в сравнении с конструкциями


первой части, схемы расширения


в [8] и в [10].


\end{notes}





Легко видеть, что имеет место цепочка равенств


$$


{(\add)^+}[{\cal L};\eta] \oplus W = {\cl}(\{{f*\eta}:f \in {B_0^+}(E,


{\cal L}) \oplus U\},{{\tau}_*}({\cal L})) =


{\cl}(\{{f*\eta}:f \in {B^+}(E,{\cal L}) \oplus V\},


{{\tau}_*}({\cal L})).


$$





С учетом положений ([6], гл.\,3) доказательство может быть сведено


к применению простейших положений из теории ТВП. Следует учесть ранее


упомянутые свойства замкнутости и, в частности, предложение 10.2,


а также свойства, отмечавшиеся в разделе 7: при


$\Om = {(\add)^+}[{\cal L};\eta] \oplus W$


имеет место вложение


${{\tau}_*}({\cal L})\mid_\Om \subset


{{\tau}_0}({\cal L})\mid_\Om$.


Наконец,


$$


{(\add)^+}[{\cal L};\eta] \oplus W = {\cl}(\{{f*\eta}:


f \in {B_0^+}(E,{\cal L}) \oplus U\},{{\tau}_\otimes}({\cal L})) =


{\cl}(\{{f*\eta}:f \in {B^+}(E,{\cal L}) \oplus V\},


{{\tau}_\otimes}({\cal L})).


$$





Доказательство практически очевидно


(см. уже отмечавшиеся свойства замкнутости множества (10.4), а также


соотношения для топологий


в ([6], c.\,45)).





\begin{props}


Имеет место ${(\add)^+}[{\cal L};\eta] \oplus W =


{\cl}(\{{f*\eta}:f \in {B_0^+}(E,{\cal L}) \oplus U\},


{\tau_{\bold B}^*}({\cal L})) =


{\cl}(\{{f*\eta}:f \in {B^+}(E,{\cal L}) \oplus V\},


{\tau_{\bold B}^*}({\cal L}))$.


\end{props}





\begin{pf}


В силу уже упоминавшегося свойства,


относящегося к вопросу об естественном погружении


(каждого из рассматриваемых


двух вариантов) пространства обычных решений в виде $*$-слабо


плотного множества в пространство ОЭ,


достаточно установить


вложение


\begin{equation}


{(\add)^+}[{\cal L};\eta] \oplus W \subset {\cl}(\{{f*\eta}:


f \in {B_0^+}(E,{\cal L}) \oplus U\},{{\tau}_{\bold B}^*}({\cal L})).


\end{equation}





Пусть $\omega$ --- элемент множества в левой части (10.12).


Поскольку множество (10.4) --- элемент семейства


${{\frak C}'_*}[{\cal L}]$, то


в обозначениях раздела 7 имеем свойство


\begin{equation}


{N_\theta}(\omega) \cap {{\bold B}_*}({\cal L}) \ne \emp,


\end{equation}


где $\theta$ есть def топология ${{\tau}_S^\otimes}({\cal L})$


при $S = {(\add)^+}[{\cal L};\eta] \oplus W$ (cм. (10.4)).


Учитывая (10.13), выберем и зафиксируем $\Theta \in {N_\theta}(\omega) \cap


{{\bold B}_*}({\cal L})$. Из простейших свойств,


упомянутых в разделе 7, имеем для некоторого множества


${{\Theta}_0} \in {N_{{\tau_*}({\cal L})}}(\omega)$ равенство


$$


\Theta = ({(\add)^+}[{\cal L};\eta] \oplus W) \cap {{\Theta}_0}.


$$


С другой стороны, по выбору $\Theta$ при некотором


${{\ld}_0} \in [0,\infty[$ для шара ${U_{{\ld}_0}}({\cal L})$


([6], c.\,45) пространства ${\Bbb A}({\cal L})$ (имеется в виду шар


в сильной норме с центром в начале координат и радиусом ${\ld}_0$)


справедливо вложение


$\Theta \subset {U_{{\ld}_0}}({\cal L})$. По известному свойству


${\tau_{\bold B}^*}({\cal L})$ ([6], c.\,54) имеем совпадение


\begin{equation}


\vt \stac {{\tau}_*}({\cal L})\mid_{{U_{{\ld}_0}}({\cal L})} =


{{\tau}_{\bold B}^*}({\cal L})\mid_{{U_{{\ld}_0}}({\cal L})}.


\end{equation}


При этом $\omega \in \Theta$ и в силу (10.14) имеет место


$$


{N_\vt}(\omega) = \{{U_{{\ld}_0}}({\cal L}) \cap H:H \in


{N_{{\tau_{\bold B}^*}({\cal L})}}(\omega)\}.


$$


Это означает, в частности, что (снова с учетом (10.14))


$\fo{H_1} \in


{N_{{\tau_{\bold B}^*}({\cal L})}}(\omega)$ $\exi{H_2} \in


{N_{{\tau_*}({\cal L})}}(\omega)$


\begin{equation}


{H_1} \cap {U_{{\ld}_0}}({\cal L}) = {H_2} \cap {U_{{\ld}_0}}({\cal L}).


\end{equation}





Далее, комбинируем (10.15) и ранее упоминавшееся


положение о плотном погружении в смысле ${{\tau}_*}({\cal L})$.


С учетом последнего и теоремы


Биркгофа ([11], гл.\,I; [20], гл.\,2; [23]) подбираем направленность


$(\Sigma,\preceq,\vp)$ в множестве ${{\Gamma}_0} \stac \{{f*\eta}:


f \in {B_0^+}(E,{\cal L}) \oplus U\}$ со свойством (сходимости по


Мору--Смиту)


\begin{equation}


(\Sigma,\preceq,\vp) @>{{\tau_*}({\cal L})}>> \omega.


\end{equation}





В (10.16) учтено свойство $\omega \in {(\add)^+}[{\cal L};\eta]


\oplus W$. Подберем ${{\sigma}_1} \in \Sigma$ из условия:


$\fo{{\sigma}'} \in \Sigma$


\begin{equation}


({{\sigma}_1} \preceq {\sigma}') \Lo


({\vp}({\sigma}') \in {{\Theta}_0}).


\end{equation}


Пусть ${\Bbb U} \in {N_{{\tau_{\bold B}^*}({\cal L})}}(\omega)$ и в


соответствии с (10.15) пусть ${\Bbb V} \in {N_{{\tau_*}({\cal L})}}(\omega)$


удовлетворяет условию


$$


{\Bbb U} \cap {U_{{\ld}_0}}({\cal L}) = {\Bbb V} \cap


{U_{{\ld}_0}}({\cal L}).


$$


Используя (10.16), подберем теперь ${{\sigma}_2} \in \Sigma$ так, что


$\fo{\wt \sigma} \in \Sigma$


\begin{equation}


({{\sigma}_2} \preceq {\wt


\sigma}) \Lo ({\vp}({\wt \sigma}) \in {\Bbb V}).


\end{equation}





Пусть ${{\sigma}_3} \in \Sigma$ есть мажоранта


$\{{{\sigma}_1};{{\sigma}_2}\}$ в $(\Sigma,\preceq)$, а


$\sigma \in \Sigma$ обладает свойством ${{\sigma}_3} \preceq \sigma$.


В силу (10.17)\; ${{\vp}(\sigma)} \in {{\Theta}_0}$ и, вместе с тем,


${{\vp}(\sigma)} \in {{\Gamma}_0}$; значит, ${{\vp}(\sigma)} \in


{(\add)^+}[{\cal L};\eta] \oplus W$.


По выбору ${\Theta}_0$ имеем ${{\vp}(\sigma)} \in \Theta$ и, как следствие,


${{\vp}(\sigma)} \in {U_{{\ld}_0}}({\cal L})$. Из (10.18) следует


${{\vp}(\sigma)} \in {\Bbb V}$ и в итоге ${{\vp}(\sigma)} \in


{\Bbb U}$. Поскольку выбор $\sigma$ был произвольным, то фактически


установлена сходимость


$$


(\Sigma,\preceq,\vp) @>{{\tau_{\bold B}^*}({\cal L})}>> \omega.


$$





По теореме Биркгофа имеем (коль скоро $(\Sigma,\preceq,\vp)$ ---


направленность в ${\Gamma}_0$) свойство


$\omega \in {\cl}({{\Gamma}_0},{\tau_{\bold B}^*}({\cal L}))$,


чем и завершается обоснование (10.12).


\end{pf}





Получили, таким образом, следующее важное положение (комбинируем


упомянутые выше свойства плотного, в смысле топологий из


${\frak M}({\cal L})$, погружения в множество (10.4)).





\begin{thms}


Имеет место следующее свойство плотности\/$:$


$\fo{\tau} \in {{\frak M}({\cal L})}$


$$


{(\add)^+}[{\cal L};\eta]


\oplus W = {\cl}(\{{f*\eta}:f \in {B_0^+}(E,{\cal L}) \oplus U\},


\tau) = {\cl}(\{{f*\eta}:f \in {B^+}(E,{\cal L}) \oplus V\},


\tau).


$$


\end{thms}





Полезно отметить достаточно очевидную теперь цепочку равенств


$$


{(\add)^+}[{\cal L};\eta] \oplus W =


{\cl}(\{{f*\eta}:f \in {B_0^+}(E,{\cal L}) \oplus U\},


{{\otimes}^{\cal L}}({\tau_{\Bbb R}})) =


{\cl}(\{{f*\eta}:f \in {B^+}(E,{\cal L}) \oplus V\},


{{\otimes}^{\cal L}}({\tau_{\Bbb R}})).


$$


Наконец,


$$


{(\add)^+}[{\cal L}] \oplus W =


{\cl}(\{{f*\eta}:f \in {B_0^+}(E,{\cal L}) \oplus U\},


{{\otimes}^{\cal L}}({\tau_\partial})) =


{\cl}(\{{f*\eta}:f \in {B^+}(E,{\cal L}) \oplus V\},


{{\otimes}^{\cal L}}({\tau_\partial})).


$$





Доказательство, по существу, сводится к применению известных


соотношений, связывающих операторы замыкания в ТП и в его п/п


(напр., [23], c.\,111). Кроме того, следует иметь в виду определения


топологий в ([6], c.\,44). В части использования двух упомянутых выше версий


тихоновского произведения полезно дополнить извеcтные ранее положения


[5],[6] подобно тому, как это было сделано в последних двух соотношениях


в части добавления к основной теореме~10.1.





В самом деле, с учетом очевидных


соотношений ([6], гл.\,3), касающихся топологий из


${\frak M}({\cal L}),\;{{\otimes}^{\cal L}}({\tau_{\Bbb R}})$


и ${{\otimes}^{\cal L}}({\tau_\partial})$, получаем:


$W$ (см. (10.3)) есть множество,


замкнутое в каждом из ТП


\begin{equation}


({{\Bbb R}^{\cal L}},{{\otimes}^{\cal L}}({\tau_{\Bbb R}})),\quad


({{\Bbb R}^{\cal L}},{{\otimes}^{\cal L}}({\tau_\partial})).


\end{equation}





Итак,


свойства плотности, упомянутые после теоремы 10.1,


характеризуют естественное погружение множеств


${B_0^+}(E,{\cal L}) \oplus U$


и ${B^+}(E,{\cal L}) \oplus V$ в ТП (10.19). Уместно дополнить


теорему 3.7.5 из [6], используя замкнутость $W$ в ТП (10.19):


\begin{multline*}


W = {\cl}(\{{f*\eta}:f \in U\},{{\otimes}^{\cal L}}({\tau_{\Bbb R}})) =


{\cl}(\{{f*\eta}:f \in V\},{{\otimes}^{\cal L}}({\tau_{\Bbb R}})) =\\


= {\cl}(\{{f*\eta}:f \in U\},{{\otimes}^{\cal L}}({\tau_\partial}))


={\cl}(\{{f*\eta}:f \in V\},{{\otimes}^{\cal L}}({\tau_\partial})).


\end{multline*}





Полученные свойства весьма универсальной плотности (см., напр., теорему


10.1) можно использовать при исследовании МП, возникающих в пространстве


ОЭ; см. в этой связи конструкции [8].


Ранее отмечалось, что множество (10.4) есть элемент семейства


${{\frak C}'_*}[{\cal L}]$ и


(см. раздел 7), как следствие,


$$


{(\add)^+}[{\cal L};\eta] \oplus W \in {\frak K}.


$$


Что же касается условия 5.1, то из теоремы 10.1 при


$$


{t_{\bold u}}={{\tau}_0}({\cal L})\mid_{{(\add)^+}[{\cal L};\eta] \oplus W}


$$


(см. (5.5)) имеем цепочку равенств


$$


{(\add)^+}[{\cal L};\eta] \oplus W = {\cl}(\{{f*\eta}:f \in


{B_0^+}(E,{\cal L}) \oplus U\},{t_{\bold u}}) =


{\cl}(\{{f*\eta}:f \in {B^+}(E,{\cal L}) \oplus V\},{t_{\bold u}}).


$$





В связи с последними свойствами уместно коснуться вопроса,


связанного с условиями плотности в смысле топологий типа


${t_{\bold u}}$. Если $({\bold T},\tau)$ есть ТП и $M \in


{{\cal P}({\bold T})}$ (т.\,е. $M$ есть п/м ${\bold T}$), то


$$


(\tau - {\dens})[M] \stac \{H \in {\cal P}({\bold T}) \mid


M \subset {\cl}(H,\tau)\}


$$


есть семейство всех п/м ${\bold T}$, всюду плотных в $M$ в смысле


$({\bold T},\tau)$.


Как и ранее, через


${\Fin}({\cal L})$ обозначаем семейство всех непустых


конечных п/м ${\cal L}$.


С учетом представления ${{\tau}_0}({\cal L})$


в ([5], c.\,81; [6], c.\,45) получаем


$\fo{A} \in {\cal P}({\Bbb A}({\cal L}))$


\begin{multline}


({{\tau}_0}({\cal L})-{\dens})[A]=\\ = \{M \in {\cal P}({\Bbb A}({\cal L}))


\mid\{(\mu\mid{\cal K}):\;\mu \in A\} \subset


\{(\mu\mid{\cal K}):\;\mu \in M\}


\ \fo{\cal K} \in {\Fin}({\cal L})\}.


\end{multline}


Из (10.20) вытекает, что $\fo{A} \in {{\cal P}({\Bbb A}({\cal L}))}$,


$\fo{M} \in ({{\tau}_0}({\cal L}) - {\dens})[A]$, $\fo{L} \in {\cal L}$


\begin{equation}


\{{\mu}(L):\mu \in A\} \subset \{{\mu}(L):\mu \in M\}.


\end{equation}





\begin{notes}


На основе (10.21) можно сделать ряд конкретных выводов об отсутствии


плотности в смысле ${{\tau}_0}({\cal L})$. Рассмотрим простой пример.


Пусть $\alpha \in [0,\infty[$ и ${\bold A} \stac \{\mu \in


{(\add)_+}[{\cal L}] \mid {{\mu}(E)} \le \alpha\}$. С учетом (10.21)


легко проверяется, что


$$


\{\mu \in {(\add)_+}[{\cal L}] \mid {{\mu}(E)} < \alpha\} \notin


({{\tau}_0}({\cal L}) - {\dens})[{\bold A}].


$$


\end{notes}





Напомним, что (см. соотношение (5.5)) версии


топологии, используемой в условии 5.1, суть топологии, индуцируемые


из $({{\Bbb R}^{\cal L}},{{\otimes}^{\cal L}}({\tau_\partial}))$


(10.19). Свойства, подобные (10.20), (10.21), легко распространяются на такие


топологии (на самом деле они справедливы для существенно более общих


конструкций, использующих тихоновские произведения экземпляров произвольного


множества, оснащаемого дискретной топологией). Сейчас отметим только, что


в обозначениях раздела 5\; $\fo{M} \in {{\cal P}'({\Bbb A}({\cal L}))}$,


$\fo{A} \in {{\cal P}(M)}$


$$


({{{\tau}_0}({\cal L})\mid_M} - \dens)[A] = \{H \in {{\cal P}(M)} \mid


\{(\mu\mid{\cal K}):\mu \in


A\} \subset\{(\mu\mid{\cal K}):\;\mu \in H\}


\ \fo{\cal K} \in {\Fin}({\cal L})\}.


$$





Наконец, подобно (10.21) имеем очевидное следствие:


$\fo{M} \in {{\cal P}'({\Bbb A}({\cal L}))}$, $\fo{A} \in {{\cal P}(M)}$,


$\fo{H} \in ({{{\tau}_0}({\cal L})\mid_M} - \dens)[A]$, $\fo{L} \in


{\cal L}$


$$


\{{\mu}(L):\mu \in A\} \subset \{{\mu}(L):\mu \in H\}.


$$








\begin{thebibliography}{99}





\bibitem{1} Варга Дж. {\em Оптимальное управление дифференциальными и


функциональными уравнениями}. -- М.: Наука, 1977. -- 624 с.





\bibitem{2} Янг Л. {\em Лекции по вариационному исчислению и теории


оптимального управления}. -- М.: Мир, 1970. -- 488 с.





\bibitem{3} Красовский Н.Н. {\em Теория управления движением}. -- М.:


Наука, 1968. -- 475 с.





\bibitem{4} Эльясберг П.Е. {\em Введение в теорию полета искусственных


спутников Земли}. -- М.: Наука, 1965. -- 540 с.





\bibitem{5} Chentsov A.G. {\em Finitely additive measures and


relaxations of extremal problems}.


- New York, London, and


Moscow: Plenum Publishing Corporation, 1996. -- 244 p.





\bibitem{6} Chentsov A.G. {\em Asymptotic attainability}. --


Dordrecht--Boston--London: Kluwer Academic Publishers, 1997. -- 322 p.





\bibitem{7}


Chentsov A.G.


{\em On approximation of asymptotic attainability domains} //


Nonsmooth and discontinuos problems of control and


optimization. -- Proc. volume from the IFAC Workshop, Chelyabinsk,


Russia, 17--20 June 1998, P.\,1--12, Pergamon.





\bibitem{8}


Chentsov A.G.


{\em Universal properties of generalized


integral constraints in the class of finitely additive measures} //


Funct. Different. Equat. -- V.\,5. -- \No\,1--2. --


P.\,69--105.





\bibitem{9} Ченцов А.Г. {\em К вопросу о корректном расширении одной задачи


о выборе плотности вероятности при ограничениях на систему математических


ожиданий} // УMH. -- 1995. -- Т.\,50. --


Bып.\,5. -- C.\,223--242.





\bibitem{10} Ченцов А.Г. {\em К вопросу о корректном расширении


некоторых неустойчивых задач управления с интегральными


ограничениями} // Изв. PАН. Сер. матем. -- 1999.


-- T.\,63. -- C.\,185--223.





\bibitem{11} Данфорд Н., Шварц


Дж.Т. {\em Линейные операторы. Общая теория}. -- М.:


Ин. лит., 1962. -- 895~с.





\bibitem{12} Самойленко А.М., Перестюк Н.А. {\em Дифференциальные уравнения


с импульсным воздействием}. -- Киев: Вища школа, 1987. -- 287 c.





\bibitem{12} Ченцов А.Г., Каширцева Т.Ю. {\em Обобщенные траектории


линейных управляемых систем с разрывными коэффициентами


при управлении} // Вестн. Челябинск. ун-та. --


1999. -- Bып.\,2. -- C.\,137--146.





\bibitem{14} Баскаков С.И. {\em Радиотехнические цепи и сигналы}.


-- М.: Высш. школа, 1988. -- 448 с.





\bibitem{15} Ченцов А.Г. {\em Конечно-аддитивные меры и задачи на минимум} //


Кибернетика. -- 1988. -- \No\,3. -- C.\,67--70.





\bibitem{16} Ченцов А.Г. {\em Двузначные меры на


полуалгебре множеств и некоторые их


приложения к бесконечномерным задачам математического


программирования} // Кибернетика. -- 1988. -- \No\,6. -- C.\,72--76.





\bibitem{17} Ченцов А.Г. {\em О представлении множеств притяжения,


возникающих при последовательном ослаблении ограничений} //


Докл. PАН. -- 1999. -- Т.\,365. -- \No\,2. --


C.\,174--177.





\bibitem{18} Ченцов А.Г. {\em Множества притяжения и их представление


в виде достижимых множеств обобщенных задач} // Изв. Уральск.


гос. ун-та. -- 1999. -- Bып.\,2. -- C.\,135--153.





\bibitem{19}


Chentsov A.G. {\em Topological constructions of extensions and


representations of attraction sets} // Proc. Steklov Inst. Math.


Suppl. Issue 1. Control Dynamic Syst., P.\,35--60.





\bibitem{20} Келли Дж.Л. {\em Общая


топология}. -- М.: Наука, 1981. -- 431 с.





\bibitem{21} Куратовский К., Мостовский А. {\em Теория множеств}. --


М.: Мир, 1970. -- 416 с.





\bibitem{22} Федорчук В.В., Филиппов В.В. {\em Общая топология. Основные


конструкции}. -- М.: Изд-во Моск. ун-та, 1988. -- 252 с.





\bibitem{23} Энгелькинг Р. {\em Общая топология}. --


М.: Мир, 1986. -- 751 с.





\bibitem{24} Ченцов А.Г. {\em Асимптотическая


достижимость при возмущении интегральных ограничений в


абстрактной задаче управления.} II // Изв. вузов. Математика. --


1995. -- \No\,3. -- С.\,62--73.





\bibitem{25} Неве Ж. {\em Математические основы теории вероятностей}. -- М.:


Мир, 1969. -- 309 с.





\bibitem{26} Понтрягин Л.С. {\em Обыкновенные


дифференциальные уравнения}. -- М.: Наука, 1965. -- 322 с.





\bibitem{27} Bhaskara Rao K.P.S., Bhaskara Rao


M. {\em Theory of charges. A study of finitely additive


measures}. -- New York: Acad. Press, 1983. -- 253 p.





\end{thebibliography}





\vskip 2cm





\post{\it Институт математики и механики}{\it Поступила}


\post{\it Уральского отделения}{29.01.2001}


\post{\it Российской Академии наук}{}





\label{end}


\end{document}


